% \iffalse meta-comment
%
% Copyright (C) 2023 by Jan Hajer
% -----------------------------------
%
% This file may be distributed and/or modified under the
% conditions of the LaTeX Project Public License, either version 1.3c
% of this license or (at your option) any later version.
% The latest version of this license is in:
%
% http://www.latex-project.org/lppl.txt
%
% and version 1.3c or later is part of all distributions of LaTeX
% version 2005/12/01 or later.
%
% \fi
%
% \iffalse
%
%<package>\NeedsTeXFormat{LaTeX2e}[2021-11-15]
%<package>\ProvidesPackage{hep-graphic}[2026/11/13 v1.6 TikZ, pgfplots and tikz-feynman setup with externalisation and figure inclusion helpers]
%<plot>\ProvidesPackage{hep-plot}[2026/11/13 v1.6 Pgfplots defaults: colour cycles, line styles and legends]
%<feynman>\ProvidesPackage{hep-feynman}[2026/11/13 v1.6 Tikz-feynman defaults]
%<documentation>\ProvidesFile{hep-graphic-documentation.tex}[2026/11/13 v1.6 HEP-graphic documentation]
%
%<*documentation>
%
\RequirePackage[l2tabu, orthodox]{nag}
\documentclass{ltxdoc}

\renewcommand*\theCodelineNo{\rmfamily\tstyle\footnotesize\arabic{CodelineNo}}
\AtBeginEnvironment{macrocode}{\renewcommand*{\ttdefault}{clmt}}
\renewcommand*{\MacroFont}{\codestyle}
\AtBeginDocument{\DeleteShortVerb{\|}}
\AtBeginDocument{\MakeShortVerb{\"}}
\EnableCrossrefs
\CodelineIndex
\RecordChanges

\usepackage{hologo}

\usepackage[parskip,oldstyle,font=10pt]{hep-paper}
\bibliography{bibliography}
%</documentation>

%<*driver>
\expandafter\newif\csname ifshort\endcsname
\shortfalse
\begin{document}
\DocInput{hep-graphic-implementation.dtx}
\end{document}
%</driver>
%
% \fi
%
% \CheckSum{288}
%
% \CharacterTable
%  {Upper-case    \A\B\C\D\E\F\G\H\I\J\K\L\M\N\O\P\Q\R\S\T\U\V\W\X\Y\Z
%   Lower-case    \a\b\c\d\e\f\g\h\i\j\k\l\m\n\o\p\q\r\s\t\u\v\w\x\y\z
%   Digits        \0\1\2\3\4\5\6\7\8\9
%   Exclamation   \!     Double quote  \"     Hash (number) \#
%   Dollar        \$     Percent       \%     Ampersand     \&
%   Acute accent  \'     Left paren    \(     Right paren   \)
%   Asterisk      \*     Plus          \+     Comma         \,
%   Minus         \-     Point         \.     Solidus       \/
%   Colon         \:     Semicolon     \;     Less than     \<
%   Equals        \=     Greater than  \>     Question mark \?
%   Commercial at \@     Left bracket  \[     Backslash     \\
%   Right bracket \]     Circumflex    \^     Underscore    \_
%   Grave accent  \`     Left brace    \{     Vertical bar  \|
%   Right brace   \}     Tilde         \~}
%
% \changes{v1.2}{2023/07/01}{First release on CTAN: hep-graphic, hep-plot, and hep-feynman with TikZ externalisation, PGF and Feynman diagram inclusion, pgfplots defaults, cycle lists, legend helpers, contour and error legend styles, and tikz-feynman settings.}
% \changes{v1.3}{2024/11/01}{Refreshed the documentation.}
% \changes{v1.4}{2025/09/01}{Applied starred command definitions.}
% \changes{v1.5}{2026/08/01}{Raised the LaTeX format requirement, expanded the documentation and README, and documented the colormap and cycle-list helpers.}
% \changes{v1.6}{2026/11/13}{Fixed the colormap positions of repeated spaced cycle lists, \cs{includepgf} without the plot option or the external library, and the spurious spaces and global plot options of \cs{includepgf} and \cs{cyclelistshift}.}
%
% \ifshort
%<*documentation>
% \fi
%
\GetFileInfo{hep-graphic.sty}

\title{The \software{hep-graphic} package\thanks{This document corresponds to \software{hep-graphic}~\fileversion.}}
\subtitle{TikZ, pgfplots and tikz-feynman setup with externalisation and figure inclusion helpers}
\author{Jan Hajer \email{jan.hajer@tecnico.ulisboa.pt}}
\date{\filedate}

% \ifshort
\begin{document}
% \fi

\newgeometry{vscale=.8, vmarginratio=3:4, includeheadfoot, left=11em, marginparwidth=4.6cm, marginparsep=3mm, right=7em}

\maketitle

\begin{abstract}
The \software{hep-graphic} package collects graphics helpers for HEP documents.
It loads \software{TikZ}, optional TikZ externalisation, and \software{tikzscale}, provides path-aware inclusion commands for PGF plots and Feynman diagrams, and can load companion \software{hep-plot} and \software{hep-feynman} modules for \software{pgfplots} defaults, colour cycles, line styles, legends, and \software{tikz-feynman} settings.
\end{abstract}

\section{Graphic}

After loading the "hep-graphic" package the \software{TikZ} \cite{pgf} package is loaded.

\DescribeMacro{extern}
The "extern="\meta{folder} option activates the TikZ "external" library and stores externalised graphics in \meta{folder}.

\DescribeMacro{plot}
\DescribeMacro{feynman}
The "plot" and "feynman" options load the \software{hep-plot} and \software{hep-feynman} companion modules.

\DescribeMacro{\graphicpath}
\DescribeMacro{\includepgf}
\DescribeMacro{\includefeynman}
The "\graphicpath"\marg{folder} macro stores the folder from which generated graphics are included.
The "\includepgf"\meta{*}\oarg{options}\marg{name} macro includes a PGF or PGFPlots graphic from that folder and passes \meta{options} to both "\pgfplotsset" and "\includegraphics"; the starred form includes the graphic without applying plotting options.
The "\includefeynman"\marg{name} macro includes a generated Feynman diagram centred for use in math displays.

\subsection{Plot}

The \software{hep-plot} package loads the \software{pgfplots} package \cite{pgfplots}, sets global plot dimensions and fonts, adjusts numeric formatting, and provides a "three panels" style for compact multi-panel figures.
It defines "rainbow" and "rainbow reversed" colormaps, spaced colour and line-style cycle-list generators, matching cycle lists for up to many curves, and the "\addlegendtitle"\marg{text} helper for legend headers.

\subsection{Feynman}

The \software{hep-feynman} package loads the \software{tikz-feynman} package \cite{tikz-feynman} and adjusts default line, momentum-arrow, and diagram styles.

% \ifshort
\printbibliography

\end{document}
%
%</documentation>
% \fi
%
% \StopEventually{
% \printbibliography
% \PrintChanges
% }
%
% \appendix
%
% \section{Implementation}
%
% \subsection{Graphic}
%
%<*package>
%
% Define a hepgraphic namespace for the options using the \software{kvoptions} package \cite{kvoptions}.
%    \begin{macrocode}
\RequirePackage{kvoptions}
\SetupKeyvalOptions{
  family=hepgraphic,
  prefix=hepgraphic@
}
%    \end{macrocode}
%
% \begin{macro}{plot}
% Define the "extern" option to activate externalisation into the declared folder.
%    \begin{macrocode}
\DeclareStringOption{extern}[.]
%    \end{macrocode}
% \end{macro}
%
% \begin{macro}{plot}
% Define the "plot" switch for loading plot code.
%    \begin{macrocode}
\DeclareBoolOption[false]{plot}
%    \end{macrocode}
% \end{macro}
%
% \begin{macro}{feynman}
% Define the "feynman" switch for loading feynman code.
%    \begin{macrocode}
\DeclareBoolOption[false]{feynman}
%    \end{macrocode}
% \end{macro}
%
%    \begin{macrocode}
\ProcessKeyvalOptions*
%    \end{macrocode}
%
% Load the \software{hep-plot} and \software{hep-feynman} packages when required.
%    \begin{macrocode}
\ifhepgraphic@plot\RequirePackage{hep-plot}\fi
\ifhepgraphic@feynman\RequirePackage{hep-feynman}\fi
%    \end{macrocode}
%
% Load the \software{tikz} package with the \software{external} library \cite{pgf}.
%    \begin{macrocode}
\RequirePackage{tikz}
\ifx\hepgraphic@extern\@empty\else
 \usetikzlibrary{external}
 \tikzexternalize[
%   optimize=false,
  only named=true,
 ]
 \tikzsetexternalprefix{\hepgraphic@extern/}
\fi
%    \end{macrocode}
%
% \begin{macro}{\graphicpath}
% Load the \software{standalone} package \cite{standalone} and define the "\graphicpath" pointing to the folder with "pgf" files.
%    \begin{macrocode}
\RequirePackage{tikzscale}
\def\hep@graphic@path{.}
\newcommand*{\graphicpath}[1]{\def\hep@graphic@path{#1}}
%    \end{macrocode}
% \end{macro}
% \begin{macro}{\includetikz}
% Define a macro to include "tikz" figures using the \software{xparse} package \cite{xparse}.
% The name of the externalised file applies only to the picture in the included "tikz" file and is cleared afterwards, so that it never names a later picture.
%    \begin{macrocode}
\RequirePackage{xparse}
\providecommand*{\tikzsetnextfilename}[1]{}
\NewDocumentCommand{\includepgf}{sO{}m}{%
 \tikzsetnextfilename{#3}%
 \IfBooleanTF{#1}{%
  \includegraphics{\hep@graphic@path/#3}%
 }{%
  \begingroup
  \ifdefined\pgfplotsset\pgfplotsset{#2}\fi
  \includegraphics[#2]{\hep@graphic@path/#3}%
  \endgroup
 }%
 \tikzsetnextfilename{}%
}
\newcommand*{\includefeynman}[1]{%
 \vcenter{\hbox{\includegraphics{\hep@graphic@path/#1}}}%
}
%    \end{macrocode}
% \end{macro}
%
%</package>
%
% \subsection{Plots}
%
%<*plot>
%
% Load the \software{pgf/tikz} package \cite{pgf}.
%    \begin{macrocode}
\RequirePackage{tikz}
%    \end{macrocode}
%
% \begin{macro}{dashdotdotdotted}
% \begin{macro}{dashdotdotdotdotted}
% Add new line styles.
%    \begin{macrocode}
% \tikzset{
%  dashdotdotdotted/.style={dash pattern=on 3pt off 2pt
%   on \the\pgflinewidth off 2pt on \the\pgflinewidth off 2pt
%   on \the\pgflinewidth off 2pt
%  },
%  dashdotdotdotdotted/.style={dash pattern=on 3pt off 2pt
%   on \the\pgflinewidth off 2pt on \the\pgflinewidth off 2pt
%   on \the\pgflinewidth off 2pt on \the\pgflinewidth off 2pt
%  },
%  dotdotdashed/.style={dash pattern=on 3pt off 2pt
%   on 3pt off 2pt on \the\pgflinewidth off 2pt
%  },
%  dotdotdotdashed/.style={dash pattern=on 3pt off 2pt
%   on 3pt off 2pt on 3pt off 2pt on \the\pgflinewidth off 2pt
%  },
% }
%    \end{macrocode}
% \end{macro}
% \end{macro}
% Change thousand separator
%    \begin{macrocode}
\pgfkeys{/pgf/number format/.cd,1000 sep={\,}}%
%    \end{macrocode}
%
% Load the \software{pgfplots} package \cite{pgfplots} and set global options.
%    \begin{macrocode}
\RequirePackage{pgfplots}
\pgfplotsset{
 compat=newest,
 width=\linewidth,
 height=\linewidth,
 enlargelimits=false,
}
%    \end{macrocode}
%
% Fix glitch.
%    \begin{macrocode}
\pgfplotsset{
 every y tick scale label/.append style={
  inner sep=1pt,
  xshift=-1pt,
  yshift=-1pt,
 },
}
%    \end{macrocode}
%
% Set default font size
%    \begin{macrocode}
\pgfplotsset{
  legend style={font=\footnotesize},
  tick label style={font=\footnotesize},
  label style={font=\small},
  title style={font=\small},
  max space between ticks=30,
}
%    \end{macrocode}
% \begin{macro}{three panels}
% Set font size three panel versions
%    \begin{macrocode}
\pgfplotsset{
  three panels/.style={
    legend style={font=\scriptsize},
    tick label style={font=\scriptsize},
    label style={font=\footnotesize},
    title style={font=\footnotesize},
    max space between ticks=25,
    /tikz/mark size=1.5pt,
    major tick length=1mm,
    minor tick length=0.66mm,
    every axis title shift=0pt,
  },
}
%    \end{macrocode}
% \end{macro}
%
% \subsubsection{Cycles}
%
% \begin{macro}{rainbow reversed}
% Define the "rainbow reversed" colormap.
% Approximate spectral colormap for visible white-light dispersion in wavelength ordering.
%    \begin{macrocode}
\usepgfplotslibrary{colormaps}
\pgfplotsset{
  colormap={rainbow reversed}{
    rgb255(410)=(38,0,79)
    rgb255(430)=(67,0,170)
    rgb255(450)=(58,0,200)
    rgb255(470)=(0,50,180)
    rgb255(490)=(0,100,120)
    rgb255(510)=(0,135,75)
    rgb255(530)=(0,150,35)
    rgb255(550)=(55,145,0)
    rgb255(570)=(174,155,0)
    rgb255(590)=(226,110,0)
    rgb255(610)=(241,0,0)
    rgb255(630)=(205,0,0)
    rgb255(650)=(141,0,0)
  },
%    \end{macrocode}
% \end{macro}
%
% \begin{macro}{rainbow}
% Define the "rainbow" colormap in frequency ordering.
%    \begin{macrocode}
  colormap={rainbow}{
    colors of colormap={
      650,630,610,590,570,550,530,510,490,470,450,430,410
      of rainbow reversed,
      source range=410:650,
      target pos={
        0,83,167,250,333,417,500,583,667,750,833,917,1000
      }
    }
  },
}
%    \end{macrocode}
% \end{macro}
%
% \begin{macro}{\l\_rainbow\_positions\_seq}
% \begin{macro}{\l\_rainbow\_target\_int}
% \begin{macro}{\rainbow\_positions\_full\_p}
% True once the current generated list has reached its requested length.
%    \begin{macrocode}
\ExplSyntaxOn
\seq_new:N \l_rainbow_positions_seq
\int_new:N \l_rainbow_target_int
\cs_new:Npn \rainbow_positions_full_p:{
  \int_compare_p:n{
    \seq_count:N \l_rainbow_positions_seq >= \l_rainbow_target_int
  }
}
%    \end{macrocode}
% \end{macro}
% \end{macro}
% \end{macro}
%
% \begin{macro}{\rainbow\_add\_position}
% Append a colormap coordinate if the requested number of colors has not yet been reached. Values are rounded only to keep printed output readable.
%    \begin{macrocode}
\cs_new_protected:Npn \rainbow_add_position:n #1{
  \int_compare:nNnT {
    \seq_count:N \l_rainbow_positions_seq } < { \l_rainbow_target_int
  }{
    \seq_put_right:Nx \l_rainbow_positions_seq{\fp_eval:n{round(#1, 6)}}
  }
}
%    \end{macrocode}
% \end{macro}
%
% \begin{macro}{\RainbowCyclePositions}
% \begin{macro}{\l\_rainbow\_denominator\_int}
% \begin{macro}{\CreateSpacedColormapCycleList}
% "\CreateSpacedColormapCycleList{<cycle list>}{<colormap>}{<count>}"
% Builds a PGFPlots cycle list by sampling <count> colors from <colormap>.
% The positions are generated by repeatedly bisecting the largest remaining
% gaps. Within each bisection level they are emitted from the outside inward:
% 0,1000,500,250,750,125,875,... This preserves even coverage while keeping
% neighboring colors visually distinct.
% Side effect: "\RainbowCyclePositions" is set to the comma-separated list of
% generated colormap coordinates for logging, inspection, or typesetting.
%    \begin{macrocode}
\newcommand{\RainbowCyclePositions}{}
\int_new:N \l_rainbow_denominator_int
\int_set:Nn \l_rainbow_denominator_int { 2 }
\NewDocumentCommand{\CreateSpacedColormapCycleList}{mmm}{
  \seq_clear:N \l_rainbow_positions_seq
  \int_set:Nn \l_rainbow_denominator_int { 2 }
  \int_set:Nn \l_rainbow_target_int { \int_max:nn { 1 } { #3 } }
  \rainbow_add_position:n { 0 }
  \rainbow_add_position:n { 1000 }
%    \end{macrocode}
% Largest-gap bisection, ordered outside-in for high adjacent contrast.
%    \begin{macrocode}
  \bool_until_do:nn { \rainbow_positions_full_p: }{
    \int_step_inline:nnnn { 1 } { 2 } { \l_rainbow_denominator_int / 2 }{
      \rainbow_add_position:n{ 1000 * ##1 / \l_rainbow_denominator_int }
      \int_compare:nNnT { ##1 } < { \l_rainbow_denominator_int - ##1 }{
        \rainbow_add_position:n{
          1000*(\l_rainbow_denominator_int-##1)/\l_rainbow_denominator_int
        }
      }
    }
    \int_set:Nn \l_rainbow_denominator_int{2*\l_rainbow_denominator_int}
  }
  \cs_gset:Npx \RainbowCyclePositions{
    \seq_use:Nn \l_rainbow_positions_seq { , }
  }
  \exp_args:Nnx \pgfplotscreateplotcyclelist { #1 }{
    [ colors~of~colormap = { \RainbowCyclePositions ~ of ~ #2 } ]
  }
}
%    \end{macrocode}
% \end{macro}
% \end{macro}
% \end{macro}
%
%
% \begin{macro}{\l\_rainbow\_line\_styles\_tl}
% \begin{macro}{\l\_rainbow\_line\_count\_int}
% \begin{macro}{\rainbow\_line\_add\_style}
% Append one line style to the line-style cycle list, stopping at the requested total count.
%    \begin{macrocode}
\tl_new:N \l_rainbow_line_styles_tl
\int_new:N \l_rainbow_line_count_int
\cs_new_protected:Npn \rainbow_line_add_style:n #1{
  \int_compare:nNnT{\l_rainbow_line_count_int}<{\l_rainbow_target_int}{
    \tl_put_right:Nn \l_rainbow_line_styles_tl { #1 \\ }
    \int_incr:N \l_rainbow_line_count_int
  }
}
%    \end{macrocode}
% \end{macro}
% \end{macro}
% \end{macro}
%
%
% \begin{macro}{\LineStyleDotLength}
% \begin{macro}{\LineStyleDashLength}
% \begin{macro}{\LineStyleGapLength}
% \begin{macro}{\l\_rainbow\_line\_pattern\_tl}
% Dot, dash, and gap lengths are functions of "\pgflinewidth", so the pattern scales when the plot line width changes.
% Filled by "\CreateSpacedColormapCycleList" with values such as 0,1000,500,250,750,...
% Append one dashed style made from "#1" dashes followed by "#2" round dots.
% Dots are zero-length dash segments with round line caps.
%    \begin{macrocode}
\newcommand{\LineStyleDotLength}{0pt}
\newcommand{\LineStyleDashLength}{4\pgflinewidth}
\newcommand{\LineStyleGapLength}{3\pgflinewidth}
\tl_new:N \l_rainbow_line_pattern_tl
\cs_new_protected:Npn \rainbow_line_add_pattern:nn #1#2{
  \tl_set:Nn \l_rainbow_line_pattern_tl { line~cap=round,dash~pattern= }
  \prg_replicate:nn { #1 }{
    \tl_put_right:Nn \l_rainbow_line_pattern_tl
    { on~\LineStyleDashLength~off~\LineStyleGapLength~ }
  }
  \prg_replicate:nn { #2 }{
    \tl_put_right:Nn \l_rainbow_line_pattern_tl
    { on~\LineStyleDotLength~off~\LineStyleGapLength~ }
  }
  \exp_args:NV \rainbow_line_add_style:n \l_rainbow_line_pattern_tl
}
%    \end{macrocode}
% \end{macro}
% \end{macro}
% \end{macro}
% \end{macro}
%
% \begin{macro}{\l\_rainbow\_line\_repeat\_int}
% \begin{macro}{\CreateLineStyleCycleList}
% "\CreateLineStyleCycleList{<cycle list>}{<count>}"
% Builds a PGFPlots cycle list of <count> line style cycle.
% Order: solid, dash, dot, dash dot, dash dash dot, dash dot dot,
% dash dash dot dot, dash dash dash dot, dash dot dot dot, \dots
% This orders styles by the smallest repetition count first, while preferring
% dash-leading patterns over dot-leading patterns.
%    \begin{macrocode}
\int_new:N \l_rainbow_line_repeat_int
\NewDocumentCommand{\CreateLineStyleCycleList}{mm}{
  \tl_clear:N \l_rainbow_line_styles_tl
  \int_zero:N \l_rainbow_line_count_int
  \int_set:Nn \l_rainbow_target_int { \int_max:nn { 1 } { #2 } }
  \rainbow_line_add_style:n { solid }
  \rainbow_line_add_pattern:nn { 1 } { 0 }
  \rainbow_line_add_pattern:nn { 0 } { 1 }
  \rainbow_line_add_pattern:nn { 1 } { 1 }
  \int_set:Nn \l_rainbow_line_repeat_int { 2 }
  \bool_until_do:nn{
    \int_compare_p:n { \l_rainbow_line_count_int >= \l_rainbow_target_int }
  }{
    \int_step_inline:nn { \l_rainbow_line_repeat_int - 1 }{
      \rainbow_line_add_pattern:nn{ \l_rainbow_line_repeat_int }{ ##1 }
      \rainbow_line_add_pattern:nn{ ##1 }{ \l_rainbow_line_repeat_int }
    }
    \rainbow_line_add_pattern:nn{
      \l_rainbow_line_repeat_int
    }{
      \l_rainbow_line_repeat_int
    }
    \int_incr:N \l_rainbow_line_repeat_int
  }
  \exp_args:NnV \pgfplotscreateplotcyclelist { #1 } \l_rainbow_line_styles_tl
}
\ExplSyntaxOff
%    \end{macrocode}
% \end{macro}
% \end{macro}
%
%
% \begin{macro}{\LineStyleCycleLength}
% \begin{macro}{line style cycle}
% \begin{macro}{\RainbowCycleLength}
% \begin{macro}{rainbow cycle}
% \begin{macro}{line styles}
% \begin{macro}{colors}
% "\RainbowCycleLength" controls how many colors are generated.
% "\LineStyleCycleLength" controls how many line style cycle are generated.
%    \begin{macrocode}
\newcommand{\LineStyleCycleLength}{64}
\CreateLineStyleCycleList{line style cycle}{\LineStyleCycleLength}
\newcommand{\RainbowCycleLength}{64}
\CreateSpacedColormapCycleList{rainbow cycle}{rainbow}{\RainbowCycleLength}
\CreateLineStyleCycleList{line styles}{\LineStyleCycleLength}
\CreateSpacedColormapCycleList{colors}{rainbow}{\RainbowCycleLength}
%    \end{macrocode}
% \end{macro}
% \end{macro}
% \end{macro}
% \end{macro}
% \end{macro}
% \end{macro}
%
%
% \begin{macro}{}
%    \begin{macrocode}
\pgfplotsset{
  cycle multiindex* list={
    rainbow cycle
    \nextlist
    line style cycle
    \nextlist
  },
}
%    \end{macrocode}
% \end{macro}
%
% \begin{macro}{colors}
% \begin{macro}{line styles}
% \begin{macro}{marks}
% \begin{macro}{horizontal marks}
% \begin{macro}{vertical marks}
% Create cycle lists
%    \begin{macrocode}
\colorlet{dark green}{green!50!black}
\colorlet{darkgreen}{dark green}
% \pgfplotscreateplotcyclelist{colors}{
%  blue, red, dark green, violet, orange, yellow!75!orange,
%  brown, black
% }
% \pgfplotscreateplotcyclelist{line styles}{
%  solid, dashed, densely dotted, dashdotted,
%  dashdotdotted, dotdotdashed, dashdotdotdotted, dotdotdotdashed, dashdotdotdotdotted
% }
\pgfplotscreateplotcyclelist{marks}{
 mark=-, mark=|, mark=Mercedes star flipped, mark=Mercedes star, mark=+,
 mark=x, mark=star, mark=asterisk, mark=10-pointed star
}
\pgfplotscreateplotcyclelist{vertical marks}{
 mark=|, mark=Mercedes star flipped, mark=Mercedes star, mark=x, mark=star,
 mark=asterisk, mark=10-pointed star
}
\pgfplotscreateplotcyclelist{horizontal marks}{
 mark=-, mark=Mercedes star flipped, mark=Mercedes star, mark=x, mark=star,
 mark=asterisk, mark=10-pointed star
}
\pgfplotscreateplotcyclelist{star marks}{
 mark=Mercedes star flipped, mark=Mercedes star, mark=x, mark=star,
 mark=asterisk, mark=10-pointed star
}
\pgfplotsset{
 cycle multiindex* list={colors\nextlist line styles},
}
%    \end{macrocode}
% \end{macro}
% \end{macro}
% \end{macro}
% \end{macro}
% \end{macro}
%
% \begin{macro}{\cyclelistshift}
% Define the "\cyclelistshift" macro, which skips one more step of the cycle list for the following plots of the current axis.
% The global counter accumulates the shifts and must be reset with "\setcounter{cyclelistshift}{0}" before the next axis.
%    \begin{macrocode}
\newcounter{cyclelistshift}
\newcommand*\cyclelistshift{%
  \addtocounter{cyclelistshift}{1}%
  \pgfplotsset{cycle list shift/.expanded=\the\value{cyclelistshift}}%
}
%    \end{macrocode}
% \end{macro}
%
% \begin{macro}{legend}
% Set the legend style.
%    \begin{macrocode}
\pgfplotsset{
 legend cell align=left,
 legend style={
  at={(1,1)},
  anchor=north east,
  inner sep=1pt,
  outer sep=6pt,
  draw=none,
  fill opacity=.9,
  draw opacity=1,
  text opacity=1,
  cells={align=left},
  /tikz/every even column/.append style={column sep=.5em},
%   fill=none,
 },
}
%    \end{macrocode}
% \end{macro}
%
% \subsubsection{Contour plots}
%
% \begin{macro}{contour legend}
% Define basic contour legend
%    \begin{macrocode}
\pgfplotsset{
 contour legend/.style={
%   contour prepared={labels=false},
  colorbar sampled line,
  colorbar style={
   mark size=7pt,
   mark options={semithick},
   tickwidth=0pt,
   subtickwidth=0pt,
   },
 },
}
%    \end{macrocode}
% \end{macro}
%
% \begin{macro}{contour legend x}
% Define horizontal contour legend.
%    \begin{macrocode}
% \usepgfplotslibrary{colormaps}
\pgfplotsset{
 contour legend x/.style={
  colorbar horizontal,
  colormap name=rainbow,
%   colormap/rainbow,
  contour legend,
  colorbar style={
   at={(0.5,1.025)},
   anchor=south,
   mark=|,
   axis x line*=top,
   axis y line=none,
   xticklabel pos=upper,
   title style={
    at={(-0.05,1)},
    anchor=east,
   },
   xlabel style={
    at={(-0.06,1)},
    anchor=south east,
   },
  },
 },
}
%    \end{macrocode}
% \end{macro}
%
%
% \begin{macro}{contour legend y}
% Define vertical contour legend.
%    \begin{macrocode}
\pgfplotsset{
 contour legend y/.style={
  contour legend,
  colorbar style={
   at={(1.025,0.5)},
   anchor=west,
   mark=-,
   axis x line=none,
   title style={
    at={(1,-0.1)},
    anchor=north west,
   },
  },
 },
}
%    \end{macrocode}
% \end{macro}
%
% \begin{macro}{contour plot x}
% \begin{macro}{contour plot y}
% Define vertical contour legend.
%    \begin{macrocode}
\pgfplotsset{
 contour plot x/.style={
  contour legend x,
  contour prepared={labels=false},
 },
 contour plot y/.style={
  contour legend y,
  contour prepared={labels=false},
 },
}
%    \end{macrocode}
% \end{macro}
% \end{macro}
%
% \begin{macro}{error legend}
% Define error legend.
%    \begin{macrocode}
\pgfplotsset{
 error legend/.style n args={3}{
  legend image code/.code={
   \draw[draw=none,fill=#1,#3](0mm,-1mm)rectangle(6mm,1mm);
   \draw[draw=#1,#2] (0mm,0mm)--(6mm,0mm);
  }
 },
}
%    \end{macrocode}
% \end{macro}
%
% \begin{macro}{\addlegendtitle}
% Define a legend title macro.
%    \begin{macrocode}
\newcommand*{\addlegendtitle}[2][]{
 \addlegendimage{empty legend}
 \addlegendentry[#1]{\hspace{-7mm}#2}
}%
%    \end{macrocode}
% \end{macro}
%
%</plot>
%
% \subsection{Feynman graphs}
%
%<*feynman>
%
% Load \software{tikz-feynman} package \cite{tikz-feynman} to enable the drawing of Feynman diagrams.
% Deactivate warning
%    \begin{macrocode}
\RequirePackage{tikz-feynman}
\tikzfeynmanset{
 compat=1.1.0,
 warn luatex=false,
}
\makeatletter\def\tikzfeynman@luatex@required@path{}\makeatother
%    \end{macrocode}
%
% Redfine the arrow style
%    \begin{macrocode}
\tikzfeynmanset{
 with arrow/.style={%
  decoration={markings,mark=at position#1with\arrow{>}},
  postaction=decorate
 },
 with reversed arrow/.style={%
  decoration={markings,mark=at position#1with\arrow{<}},
  postaction=decorate
 },
 momentum/arrow style={->},
}
%    \end{macrocode}
%
%</feynman>
%
% \section{Tests}
%
%<*test>
%
%    \begin{macrocode}
\documentclass{article}

\usepackage{hep-graphic}

\begin{document}

\end{document}
%    \end{macrocode}
%
%</test>
%
% \section{Readme}
%
%<*readme>
%
%    \begin{macrocode}
# The `hep-graphic` package

TikZ, pgfplots and tikz-feynman setup with externalisation and figure
inclusion helpers

## Introduction

The `hep-graphic` package collects graphics helpers for HEP documents. It
loads TikZ, optional TikZ externalisation, and `tikzscale`, provides
path-aware inclusion commands for generated PGF plots and Feynman
diagrams, and can load companion `hep-plot` and `hep-feynman` modules for
`pgfplots` defaults, colour cycles, line styles, legends, and
`tikz-feynman` settings.

It can be loaded with `\usepackage{hep-graphic}`. The `plot` option loads
the plot module, the `feynman` option loads the Feynman-diagram module,
and `extern=<folder>` activates TikZ externalisation into the selected
folder.

## Author

Jan Hajer

## License

This file may be distributed and/or modified under the conditions of the
`LaTeX` Project Public License, either version 1.3c of this license or
(at your option) any later version. The latest version of this license is
in `http://www.latex-project.org/lppl.txt` and version 1.3c or later is
part of all distributions of LaTeX version 2005/12/01 or later.
%    \end{macrocode}
%
%</readme>
%
% \Finale

\endinput

% \PrintIndex
% makeindex -s gglo.ist -o hep-graphic-implementation.gls hep-graphic-implementation.glo
% makeindex -s gglo.ist -o hep-graphic-implementation.ind hep-graphic-implementation.idx
